\documentclass[11pt,a4paper]{article}
\usepackage{isabelle,isabellesym}
\usepackage{amsfonts, amsmath, amssymb}

% this should be the last package used
\usepackage{pdfsetup}
\usepackage[shortcuts]{extdash}

% urls in roman style, theory text in math-similar italics
\urlstyle{rm}
\isabellestyle{rm}


\begin{document}

\title{Fast Chinese Remaindering via Product Trees}
\author{Manuel Eberl}
\maketitle

\begin{abstract}
Product trees are a simple data structure that, broadly speaking, allows dealing with the product
of a large number of (typically coprime) integers, or more generally elements of a Euclidean ring.

This is particularly useful for efficient simultaneous 
modular reduction, i.e.\ computing $x\mathbin{\text{mod}} m_i$ for a large number of moduli $m_i$ at the
same time for a fixed $x$, and the inverse operation thereof, i.e.\ modular reconstruction 
(also known as ``Chinese Remaindering'').

The algorithms formalised are adapted from the book ``Modern Computer Algebra'' by von zur
Gathen and Gerhard~\cite[Chapter 10]{gathen}.
\end{abstract}

\tableofcontents

\newpage
\parindent 0pt\parskip 0.5ex

\input{session}

\raggedright
\bibliographystyle{abbrv}
\bibliography{root}

\end{document}

